\documentclass[11pt]{article}

\usepackage[a4paper,margin=1in]{geometry}
\usepackage{amsmath,amssymb,amsthm,mathtools}
\usepackage[T1]{fontenc}
\usepackage{lmodern}
\usepackage{microtype}
\usepackage{hyperref}

\newtheorem{lemma}{Lemma}
\newtheorem{proposition}{Proposition}

\DeclareMathOperator{\tridiag}{tridiag}

\title{Monotonicity of the Largest Real-Axis Barrier}
\author{}
\date{}

\begin{document}
\maketitle

Fix \(0<a<1\), and put
\[
h=1+a,
\qquad
\alpha=a+a^{-1}.
\]
We prove that
\[
\tau_{n+1}(a)\leq \tau_n(a)
\qquad(n\geq2).
\]

\section*{1. A symmetric signed singular-value problem}

Let \(J_n\) denote the reversal matrix,
\[
J_ne_j=e_{n+1-j}.
\]
For
\[
B_n(x)=xI_n-A_n(a)
\]
define
\[
C_n(x)=B_n(x)J_n.
\]
Since
\[
J_nA_n(a)^TJ_n=A_n(a),
\]
the matrix \(C_n(x)\) is real symmetric. Moreover,
\[
B_n(x)^TB_n(x)
 =J_nC_n(x)^2J_n.
\]
Consequently, the singular values of \(B_n(x)\) are the absolute values of the
eigenvalues of \(C_n(x)\).

Introduce
\[
F_n(x,t)=\det\bigl(xI_n-A_n(a)-tJ_n\bigr).
\]
Because \(C_n(x)-tI_n=(xI_n-A_n(a)-tJ_n)J_n\), the zeros in \(t\) of
\(F_n(x,t)\) are, up to the harmless constant \(\det J_n\), the eigenvalues of
\(C_n(x)\). In particular,
\[
\det\bigl(B_n(x)^TB_n(x)-t^2I_n\bigr)
 =F_n(x,t)F_n(x,-t).
\tag{1}
\]

We shall use the sign of the eigenvalue of \(C_n(x)\) that vanishes at the two
ends of a given spectral gap. To determine it, let
\[
D_n=\operatorname{diag}(1,\sqrt a,a,\ldots,a^{(n-1)/2})
\]
and
\[
S_n=\tridiag(\sqrt a,0,\sqrt a).
\]
Then
\[
A_n(a)=D_nS_nD_n^{-1}
\]
and
\[
C_n(x)
 =a^{-(n-1)/2}D_n(xI_n-S_n)J_nD_n.
\tag{2}
\]
Thus \(C_n(x)\) is congruent, up to a positive scalar, to the commuting
symmetric product \((xI_n-S_n)J_n\).

The eigenvector of \(S_n\) associated with the \(k\)-th eigenvalue in increasing
order has reversal parity \((-1)^{n-k}\). Hence, on the gap
\((\lambda_{n,k},\lambda_{n,k+1})\), the two diagonal entries of
\((xI_n-S_n)J_n\) which change sign at the endpoints have the common sign
\[
\sigma_{n,k}=(-1)^{n-k}=(-1)^{n+k}.
\tag{3}
\]
The corresponding eigenvalue of \(C_n(x)\) is therefore
\(\sigma_{n,k}\mu_n(x;a)\).

We next reduce its level curves to a scalar oscillation problem.

\section*{2. Evaluation of the determinant}

Set
\[
c=x^2+(1-a)^2-t^2,
\qquad
d=2a^2+2ac-h^2x^2.
\]
Define \(R_j=0\) for \(j<1\), \(R_1=1\), and, for \(m\geq2\),
\[
R_m
 =cR_{m-1}+dR_{m-2}+a^2cR_{m-3}-a^4R_{m-4}.
\tag{4}
\]
A continuant expansion after pairing the first and last coordinates, the
second and penultimate coordinates, and so forth, gives
\begin{align}
F_{2m}(x,t)
 &=R_{m+1}-\bigl(1-a+a^2+ht\bigr)R_m
   -a\bigl(1-a+a^2+ht\bigr)R_{m-1}+a^3R_{m-2},
\tag{5}\\
F_{2m+1}(x,t)
 &=(x-t)R_{m+1}
   -\bigl(x(1+a^2)+2at\bigr)R_m
   +a^2(x-t)R_{m-1}.
\tag{6}
\end{align}
For completeness, these identities may also be verified directly from the
generating functions
\begin{align}
\sum_{m\geq0}F_{2m}(x,t)z^m
 &=
 \frac{1-\bigl(1-a+a^2+ht\bigr)z
       -a\bigl(1-a+a^2+ht\bigr)z^2+a^3z^3}
 {1-cz-dz^2-a^2cz^3+a^4z^4},
\tag{7}\\
\sum_{m\geq0}F_{2m+1}(x,t)z^m
 &=
 \frac{(x-t)-\bigl(x(1+a^2)+2at\bigr)z+a^2(x-t)z^2}
 {1-cz-dz^2-a^2cz^3+a^4z^4}.
\tag{8}
\end{align}
Indeed, expanding the two right-hand sides and comparing coefficients gives
the folded continuant recurrence and the correct initial determinants
\(F_0=1\), \(F_1=x-t\), \(F_2=x^2-a-t^2-2t\).

Let \(u,v\) be the roots of
\[
z^2-\frac ca z-\left(\frac d{a^2}+2\right)=0,
\qquad u\geq v.
\tag{9}
\]
The denominator in (7)--(8) factors accordingly. If \(U_j\) denotes the
Chebyshev polynomial of the second kind, normalized by
\[
U_{-1}=0,\qquad U_0=1,\qquad
U_{j+1}(z/2)=zU_j(z/2)-U_{j-1}(z/2),
\]
then
\[
R_m
 =a^{m-1}
 \frac{U_m(u/2)-U_m(v/2)}{u-v}.
\tag{10}
\]
Put
\[
\Phi_m(z)=U_m(z/2)+U_{m-1}(z/2).
\]
Substitution of (10) into (5)--(6), followed only by the Chebyshev recurrence,
gives
\begin{align}
F_{2m}(x,t)
 &=
 \frac{a^{m-1}}{u-v}
 \Bigl(
  [a(u-\alpha)-ht]\Phi_m(u)
  -[a(v-\alpha)-ht]\Phi_m(v)
 \Bigr),
\tag{11}\\
F_{2m+1}(x,t)
 &=
 \frac{a^m}{u-v}
 \Bigl(
  [x(u-\alpha)-t(u+2)]U_m(u/2)
 \notag\\[-2mm]
 &\hspace{42mm}
  -[x(v-\alpha)-t(v+2)]U_m(v/2)
 \Bigr).
\tag{12}
\end{align}

Define
\[
r(z)=\frac{a(\alpha-z)}h,
\qquad
p(z)=\frac{a(z+2)}h=h-r(z).
\tag{13}
\]
Using (9), one obtains
\[
x^2=p(u)p(v),
\qquad
t^2=r(u)r(v).
\tag{14}
\]
On the branch corresponding to the negative half of the real axis, write
\[
r_1=r(u),\qquad r_2=r(v),
\qquad 0\leq r_1\leq r_2\leq h.
\]
Then
\[
x=-\sqrt{(h-r_1)(h-r_2)},
\qquad
t=\sqrt{r_1r_2}.
\tag{15}
\]

For \(z(r)=\alpha-hr/a\), define
\begin{align}
Y_{2m}(r)
 &=\sqrt r\,
   \bigl(U_m(z(r)/2)+U_{m-1}(z(r)/2)\bigr),
\tag{16}\\
Y_{2m+1}(r)
 &=\sqrt{r(h-r)}\,U_m(z(r)/2).
\tag{17}
\end{align}
Equations (11)--(15) show that
\[
F_n\bigl(x,\sigma_{n,k}t\bigr)=0
\quad\Longleftrightarrow\quad
Y_n(r_2)=(-1)^kY_n(r_1).
\tag{18}
\]

The zeros relevant to the negative eigenvalues are
\[
r_{n,j}
 =\frac{1+a^2-2a\cos\frac{2j\pi}{n+1}}h,
\qquad
1\leq j\leq \left\lfloor\frac n2\right\rfloor.
\tag{19}
\]
Indeed, if \(r_1=0\) and \(r_2=r_{n,j}\), then (15) gives
\[
x=-2\sqrt a\cos\frac{j\pi}{n+1}.
\]
These are precisely the negative eigenvalues of \(A_n(a)\).

The Chebyshev identity
\[
U_j(s)^2-U_{j-1}(s)U_{j+1}(s)=1
\tag{20}
\]
implies
\begin{align}
\frac haY_{2m-1}(r)Y_{2m+1}(r)&=Y_{2m}(r)^2-r,
\tag{21}\\
Y_{2m}(r)Y_{2m+2}(r)&=\frac haY_{2m+1}(r)^2-r.
\tag{22}
\end{align}
Thus the zeros of consecutive functions strictly interlace, and between two
successive zeros the absolute value has one and only one maximum. It follows
from (18), first locally at \(t=0\) and then by continuation using
(21)--(22), that the signed eigenvalue specified in (3) is represented by the
unique component of
\[
Y_n(r_2)=(-1)^kY_n(r_1)
\tag{23}
\]
issuing from
\[
(0,r_{n,k})
\quad\text{and}\quad
(0,r_{n,k+1}).
\]
For the last negative gap, the second endpoint is interpreted as \(r_2=h\)
when \(n\) is odd. When \(n\) is even, the central half-gap terminates on
\(r_2=h\). Hence
\[
\beta_{n,k}(a)
 =\max_{\Gamma_{n,k}}\sqrt{r_1r_2},
\tag{24}
\]
where \(\Gamma_{n,k}\) is the corresponding component of (23). The
interlacing in (21)--(22) also shows that no other component can represent the
least singular value: every other component lies between two additional
zeros and therefore belongs to a more distant eigenvalue of \(C_n(x)\).

\section*{3. The scalar oscillation}

Put
\[
\kappa=\frac{(1-a)^2}{4a}>0
\]
and parametrize
\[
r(\theta)=\frac{4a}{h}\bigl(\kappa+\sin^2\theta\bigr),
\qquad 0\leq\theta\leq\frac{\pi}{2}.
\tag{25}
\]
For \(0\leq r<(1-a)^2/h\), the same formula is used with
\(\theta=i\eta\). Since
\[
z(r(\theta))=2\cos(2\theta),
\]
the identities
\[
U_m(\cos2\theta)+U_{m-1}(\cos2\theta)
 =\frac{\sin((2m+1)\theta)}{\sin\theta},
\]
and
\[
\sin(2\theta)U_m(\cos2\theta)=\sin((2m+2)\theta)
\]
show that, after multiplication by a positive constant depending only on the
parity of \(n\), both (16) and (17) reduce to
\[
G_M(\theta)=w(\theta)\sin(M\theta),
\qquad
M=n+1,
\tag{26}
\]
where
\[
w(\theta)=\sqrt{1+\kappa\csc^2\theta}.
\tag{27}
\]
On the hyperbolic part this means
\[
G_M(i\eta)
 =w_h(\eta)\sinh(M\eta),
\qquad
w_h(\eta)=\sqrt{\kappa\operatorname{csch}^2\eta-1}.
\tag{28}
\]
Only absolute values matter in the following argument.

We now prove the comparison of the principal equal-amplitude components for
successive frequencies.

\begin{lemma}[Nonterminal lobes]
Let \(M\geq3\), and consider a principal equal-amplitude pair for \(G_M\)
whose right coordinate lies in a lobe strictly before the terminal central
lobe. There is an equal-amplitude pair for \(G_{M-1}\) whose two \(r\)-coordinates
are not smaller. Consequently, its product \(r_1r_2\) is not smaller.
\end{lemma}

\begin{proof}
Put
\[
c=\frac{M}{M-1}>1.
\]
First suppose that both coordinates are trigonometric and are represented by
\(0<\theta_1<\theta_2<\pi/(2c)\). Map them to
\[
\theta'_j=c\theta_j.
\]
Then
\[
\sin((M-1)\theta'_j)=\sin(M\theta_j).
\]
Let
\[
Q_c(\theta)=\frac{w(c\theta)}{w(\theta)}.
\]
We claim that \(Q_c\) is strictly increasing on
\((0,\pi/(2c))\). Indeed,
\[
\frac{d}{d\theta}\log Q_c(\theta)^2
 =2\bigl(\delta(\theta)-c\delta(c\theta)\bigr),
\qquad
\delta(u)=\frac{\kappa\cot u}{\kappa+\sin^2u}.
\tag{29}
\]
Now
\[
u\delta(u)
 =\frac{u\cot u}{\kappa+\sin^2u}
\]
is strictly decreasing: \(u\cot u\) is strictly decreasing, while
\(\kappa+\sin^2u\) is strictly increasing. Hence
\(\delta(\theta)>c\delta(c\theta)\), proving the claim.

Because \(\theta_2>\theta_1\), the mapped right amplitude for \(G_{M-1}\) is
larger than the mapped left amplitude. Moving the right coordinate from
\(c\theta_2\) toward the right zero of its lobe decreases that amplitude
continuously to zero. Therefore it assumes the left amplitude at some point
\(\widehat\theta_2\geq c\theta_2\). Since \(r(\theta)\) is increasing, both
new \(r\)-coordinates are at least the original ones.

If the first coordinate is hyperbolic, keep its \(r\)-value fixed. On that
coordinate the ratio of the \((M-1)\)-amplitude to the \(M\)-amplitude is
\[
\frac{\sinh((M-1)\eta)}{\sinh(M\eta)}
 <\frac{M-1}{M}.
\tag{30}
\]
On the mapped trigonometric right coordinate the corresponding ratio is
\(Q_c(\theta_2)\). Since
\[
\lim_{\theta\downarrow0}Q_c(\theta)=\frac1c=\frac{M-1}{M}
\]
and \(Q_c\) is increasing, the right amplitude again exceeds the left one.
The same intermediate-value argument completes the proof.
\end{proof}

It remains to compare the two possible terminal configurations.

\begin{lemma}[Odd terminal lobe versus even central half-gap]
Suppose \(M\) is even. The product maximum for the terminal lobe of \(G_M\)
is no larger than the product maximum for the central half-lobe of
\(G_{M-1}\).
\end{lemma}

\begin{proof}
Put
\[
\theta_c=\frac{(M-1)\pi}{2M}.
\]
If the right coordinate satisfies \(\theta_2\leq\theta_c\), the preceding
lemma applies. It remains to consider \(\theta_2\geq\theta_c\).

We need a bound for the amplitude on the first inverse branch. Put
\[
R_M=
\frac{w\bigl(\pi/(M-1)\bigr)}{w(\pi/M)}.
\tag{31}
\]
We claim that, up to and including the first maximum of \(G_M\),
\[
\frac{G_{M-1}(u)}{G_M(u)}\leq R_M.
\tag{32}
\]
On the hyperbolic part, the left side is at most \((M-1)/M\), while
\[
R_M\geq
\frac{\sin(\pi/M)}{\sin(\pi/(M-1))}
\geq\frac{M-1}{M}.
\tag{33}
\]
On the trigonometric part,
\[
\frac{G_{M-1}(u)}{G_M(u)}
 =\frac{\sin((M-1)u)}{\sin(Mu)}
\]
is increasing, so it is enough to check the first maximum of \(G_M\). At that
point set
\[
A=\sin^2\frac{\pi}{M},
\qquad
B=\sin^2\frac{\pi}{M-1},
\qquad
E=\frac{\kappa}{\kappa+\sin^2u}.
\]
The stationary equation for \(G_M\) gives
\[
E=\frac{M\tan u}{\tan(Mu)}
\]
and hence
\[
\frac{\sin((M-1)u)}{\sin(Mu)}
 =\cos u\left(1-\frac EM\right).
\]
Therefore
\[
1-
\left(\frac{\sin((M-1)u)}{\sin(Mu)}\right)^2
 \geq \frac{E(2M-1)}{M^2}.
\tag{34}
\]
On the other hand,
\[
R_M^2=\frac{1+\kappa/B}{1+\kappa/A}
\]
and, because \(A\geq\sin^2u\),
\[
1-R_M^2
 =\left(1-\frac AB\right)\frac{\kappa}{A+\kappa}
 \leq E\left(1-\frac AB\right)
 \leq\frac{E(2M-1)}{M^2}.
\tag{35}
\]
Here the last inequality follows from
\[
\frac{\sin(\pi/M)}{\sin(\pi/(M-1))}
 \geq\frac{M-1}{M}.
\]
Equations (34)--(35) prove (32).

The monotonicity of \(Q_c\), applied between \(\pi/M\) and \(\theta_c\),
also gives
\[
R_M\leq\frac{w(\pi/2)}{w(\theta_c)}.
\tag{36}
\]
The amplitude of the terminal \(M\)-lobe for
\(\theta_2\geq\theta_c\) is at most \(w(\theta_c)\). Hence (32)--(36) imply
that the \((M-1)\)-amplitude at the same left \(r\)-coordinate is at most
\(w(\pi/2)\), which is the amplitude at the central boundary \(r=h\).
Thus the central equal-amplitude pair for \(G_{M-1}\) has left coordinate at
least as large and right coordinate equal to \(h\). Its product is therefore
at least the original product.

Finally, the maximum on the even central half-lobe really occurs at its
boundary. Since \(M-1\) is odd, write
\[
\theta=\frac{\pi}{2}-v,
\qquad
0\leq v\leq\frac{\pi}{2(M-1)}.
\]
Then
\[
w(\theta)^2\sin^2((M-1)\theta)
 =\left(1+\frac{\kappa}{\cos^2v}\right)
   \cos^2((M-1)v)
 \leq 1+\kappa=w(\pi/2)^2,
\]
because \(\cos((M-1)v)\leq\cos v\). This completes the proof.
\end{proof}

\begin{lemma}[Even central half-gap versus odd terminal lobe]
Let \(L\geq2\) be even. The product maximum for the central half-lobe of
frequency \(L+1\) is no larger than the product maximum for the terminal lobe
of frequency \(L\).
\end{lemma}

\begin{proof}
It is convenient to write
\[
C_j(u)=\frac{\sin(ju)}{\sin u},
\]
with the usual hyperbolic continuation when \(u=i\eta\), and
\[
\rho(u)=\frac{r(u)}h
       =\frac{\kappa+\sin^2u}{\kappa+1}.
\tag{37}
\]
Let the maximizing central pair for frequency \(L+1\) be \((q,h)\), and let
\(\zeta\) be the first-sheet parameter corresponding to \(q\). Equality of
the two amplitudes says
\[
\rho(\zeta)C_{L+1}(\zeta)^2=1.
\tag{38}
\]

Set
\[
\beta=\frac{\pi}{2}-\frac{\pi}{2L},
\qquad
B=r(\beta),
\qquad
b=\frac Bh=\rho(\beta),
\qquad
A=\frac{qh}{B}=\frac qb.
\tag{39}
\]
Thus \(AB=qh\). Let \(\eta\) be the first-sheet parameter corresponding to
\(A\), and define \(\eta_0\) by
\[
\sin^2\eta_0=\frac{\sin^2\zeta}{\sin^2\beta},
\tag{40}
\]
again with hyperbolic continuation when necessary. Direct substitution of
(37)--(40) gives
\[
\sin^2\eta-\sin^2\eta_0
 =
 \frac{\kappa\cos^2\beta
       \bigl(\sin^2\beta-\sin^2\zeta\bigr)}
      {\sin^2\beta\bigl(\kappa+\sin^2\beta\bigr)}
 \geq0.
\tag{41}
\]
On the first sheet \(C_L\) decreases as \(r\) increases. Indeed, before its
first zero the logarithmic derivative of \(\sin(Lu)/\sin u\) is
\(L\cot(Lu)-\cot u<0\), and the hyperbolic statement follows in the same way.
Thus
\[
C_L(\eta)\leq C_L(\eta_0).
\tag{42}
\]

We claim that
\[
\frac{C_L(\eta_0)}{C_{L+1}(\zeta)}
 \leq\sin\beta.
\tag{43}
\]
In the trigonometric case, (38) places
\(0\leq\zeta\leq\pi/(2(L+1))\). It is enough to show
\[
L\eta_0\leq(L+1)\zeta.
\]
By (40), this follows from
\[
\frac{\sin((L+1)\zeta/L)}{\sin\zeta}
 \geq\frac1{\sin\beta}.
\tag{44}
\]
The left side decreases in \(\zeta\), so it suffices to check the right
endpoint. There (44) becomes
\[
\frac{\sin(\pi/(2L))}{\sin(\pi/(2(L+1)))}
 \geq\frac1{\cos(\pi/(2L))},
\]
or equivalently
\[
\frac12\sin\frac{\pi}{L}
 \geq\sin\frac{\pi}{2(L+1)}.
\tag{45}
\]
For \(L=2\) this is an equality, and for \(L>2\) it follows from the strict
concavity of \(\sin\) on \([0,\pi/2]\).

If \(\zeta=i\chi\), then
\[
\frac{\sinh((L+1)\chi/L)}{\sinh\chi}
 \geq\frac{L+1}{L}
 \geq\frac1{\cos(\pi/(2L))},
\]
which proves the hyperbolic version of (43).

At \(B\), since \(|\sin(L\beta)|=1\), the terminal \(L\)-amplitude is
\[
\frac{\sqrt B}{\sin\beta}.
\]
At \(A\), it is \(\sqrt A\,C_L(\eta)\). From (38)--(39),
\[
\sqrt{\frac AB}
 =\frac1{bC_{L+1}(\zeta)}.
\]
Using (42), (43), and \(b\geq\sin^2\beta\), we obtain
\[
\frac{\sqrt A\,C_L(\eta)}
     {\sqrt B/\sin\beta}
 \leq
 \frac{\sin\beta}{b}
 \frac{C_L(\eta_0)}{C_{L+1}(\zeta)}
 \leq\frac{\sin^2\beta}{b}
 \leq1.
\tag{46}
\]
Thus, on the fixed-product curve \(r_1r_2=qh\), at \(r_2=B\) the right
terminal amplitude is at least the left amplitude. At \(r_2=h\), the right
terminal amplitude is zero while the left amplitude is positive.
Continuity therefore gives an equal-amplitude terminal pair on the same
fixed-product curve. Its product is exactly \(qh\), proving the desired
comparison.
\end{proof}

\section*{4. Completion of the proof}

The relation
\[
J_n(xI_n-A_n(a))J_n=xI_n-A_n(a)^T
\]
and invariance of singular values under transposition show that
\[
\mu_n(-x;a)=\mu_n(x;a).
\tag{47}
\]
Therefore
\[
\beta_{n,k}(a)=\beta_{n,n-k}(a).
\tag{48}
\]

Let \(N=n+1\).

If \(N=2m+1\) is odd, every nonterminal negative gap of size \(N\) is
dominated, by the first lemma, by the corresponding gap of size \(N-1=2m\).
The last negative gap is dominated, by the second lemma, by the central gap
of size \(2m\). Hence every negative-half barrier for size \(N\) is at most
\(\tau_{N-1}(a)\).

If \(N=2m+2\) is even, the noncentral negative gaps are again covered by the
first lemma. The central gap is dominated, by the third lemma, by the last
negative gap of size \(N-1=2m+1\). Hence every negative or central barrier for
size \(N\) is at most \(\tau_{N-1}(a)\).

By the symmetry (48), this covers all gaps. Therefore
\[
\beta_{N,k}(a)\leq\tau_{N-1}(a)
\qquad(1\leq k\leq N-1).
\]
Taking the maximum over \(k\) gives
\[
\tau_N(a)\leq\tau_{N-1}(a).
\]
Replacing \(N\) by \(n+1\), we conclude that
\[
\boxed{\tau_{n+1}(a)\leq\tau_n(a)}
\qquad
(0<a<1,\ n\geq2).
\]

\end{document}
